\documentclass[11pt]{article}
\usepackage{notes} % place notes.sty in the same folder
\definecolor{AccN}{HTML}{0D5C6F}
\fancyhead[L]{\small\color{NavyN}\nouppercase{\leftmark}}\fancyhead[R]{\small\color{GreyN}Instructor Manual --- Lecture 4}
\usepackage{longtable,multicol,needspace}
\setlength{\emergencystretch}{2em}
\newtcolorbox{instr}[1]{nb,colback=AccN!6,colframe=AccN,title={#1}}
\newtcolorbox{prob}{nb,colback=BluN!5,colframe=BluN,title={Problem statement}}
\newtcolorbox{soln}{nb,colback=GrnN!5,colframe=GrnN,title={Model solution and marking points}}
\newtcolorbox{sayit}{nb,colback=white,colframe=NavyN,title={Read aloud to students}}
\newtcolorbox{live}[1]{nb,colback=AccN!4,colframe=AccN,title={#1}}
\begin{document}
\begin{titlepage}\centering\vspace*{2.2cm}
{\color{AccN}\large\bfseries LECTURE 4\par}\vspace{0.6cm}
{\color{NavyN}\Huge\bfseries Unsupervised Learning Part 1:\\Clustering \& Dimensionality Reduction\par}\vspace{0.6cm}
{\color{AccN}\rule{0.7\textwidth}{1.5pt}}\par\vspace{0.6cm}
{\LARGE Instructor Manual\par}\vspace{0.4cm}
{\large\begin{minipage}{0.85\textwidth}\centering k-Means \textbullet\ DBSCAN \textbullet\ GMM \textbullet\ Hierarchical \textbullet\ PCA \textbullet\ t-SNE \textbullet\ UMAP\\Case study: OPFOR tactical archetype discovery, N = 8,000 trajectories, JANUS wargame simulator, 12 features\end{minipage}\par}\vspace{2cm}
{\large AK DHAMIJA, Researcher \& Instructor, AI \& Machine Learning\par}\vspace{0.2cm}
{\normalsize \href{mailto:dhamija.ak@gmail.com}{dhamija.ak@gmail.com} \quad\textbullet\quad 9811453163\par}\vspace{0.3cm}
{Unit 4: Simulation Strategies in Warfare --- AI \& Machine Learning in Warfare\par}\vfill
\begin{tcolorbox}[colback=NavyN!4,colframe=NavyN,boxrule=0.6pt,arc=2pt,width=0.9\textwidth]\small
\textbf{This manual contains:} the 45-minute session plan with timings; a slide-by-slide teaching guide for all 26 slides; key messages, facilitation tips, anticipated questions with model answers, misconceptions and in-class activities; live demonstrations with the interactive studio; consistency notes on the source material; the complete quiz answer key with rationale and scoring rules; and the assignment with marking scheme and full model solutions.
\end{tcolorbox}
\end{titlepage}
\tableofcontents\newpage


\needspace{12\baselineskip}\section{Lecture at a Glance}

\begin{center}\small\renewcommand{\arraystretch}{1.25}\begin{tabularx}{\textwidth}{@{}>{\bfseries}p{3.3cm}>{\raggedright\arraybackslash}X@{}}\toprule

\textcolor{NavyN}{Lecture}&4 --- Unsupervised Learning Part 1: Clustering \& Dimensionality Reduction\\
\textcolor{NavyN}{Duration}&45 minutes: 5 min opening Q\&A, 30 min teaching, 10 min auto-quiz\\
\textcolor{NavyN}{Slides}&26 (Beamer PDF, 16:9), 10 sections\\
\textcolor{NavyN}{Prerequisite}&Lecture~3; students should have attempted its quiz and assignment.\\
\textcolor{NavyN}{Student materials}&\texttt{Lecture4\_Slides\_Beamer.pdf}; \texttt{Lecture4\_Reference\_Text.pdf}; \texttt{Lecture4\_Quiz.html}; \texttt{Lecture4\_Answers\_Assignment.html} (release after the quiz)\\
\textcolor{NavyN}{Interactive studio}&\texttt{Lecture4\_Clustering\_Lab.html} (live demonstrations in Section~\ref{sec:live})\\
\textcolor{NavyN}{Assessment}&Auto-quiz: 15 MCQs, 40~s each, $+1/-0.25$. Assignment: 5 problems, 100 marks, plus a 20-mark bonus problem.\\\bottomrule\end{tabularx}\end{center}

\subsection*{Overview}Lecture~4 turns to \emph{unsupervised} learning: finding structure when nobody has labelled the data. The running case clusters 8,000 OPFOR manoeuvre trajectories from JANUS simulations described by 12 features. $k$-Means (with the elbow and silhouette) finds five archetypes --- including INFILTRATION, a tactic that was not in the analysts' label scheme; DBSCAN isolates 142 deception manoeuvres as noise; GMMs give soft memberships and model selection by AIC/BIC; Ward's dendrogram agrees with $k$-Means (ARI 0.847). PCA, t-SNE and UMAP complete the toolkit for seeing high-dimensional data. The lecture's discipline is validation: clusters are hypotheses until confirmed by several methods and by domain experts.

\subsection*{Learning outcomes}By the end of the lecture students should be able to:\begin{enumerate}\item Run Lloyd's algorithm by hand and explain why it converges.\item Choose $k$ with the elbow and silhouette, and interpret silhouette values.\item Explain DBSCAN core, border and noise points and tune $\varepsilon$ and MinPts.\item Describe the EM algorithm for GMMs and select components with AIC/BIC.\item Compute and interpret PCA explained variance; use t-SNE/UMAP responsibly.\item Validate clusters using internal indices, ARI and expert review.\end{enumerate}

\subsection*{Syllabus mapping}\begin{itemize}\item Unsupervised learning and clustering --- \textbf{pattern discovery} (OPFOR archetype discovery, \textbf{unusual force behaviour identification} via DBSCAN noise)\end{itemize}

\needspace{12\baselineskip}\section{Session Plan (45 Minutes)}

\begin{center}\begin{tikzpicture}[x=0.3444cm,y=1cm]
\fill[AmbN!40](0,0)rectangle(5,1);\node[font=\scriptsize\bfseries,align=center]at(2.5,0.5){Q\&A\\recap};
\fill[AccN!25](5.00,0)rectangle(7.50,1);\draw[white,thick](5.00,0)--(5.00,1);
\node[font=\tiny,rotate=90,anchor=center]at(6.25,0.5){\parbox{0.9cm}{\centering \S1}};
\fill[NavyN!18](7.50,0)rectangle(14.50,1);\draw[white,thick](7.50,0)--(7.50,1);
\node[font=\tiny,rotate=90,anchor=center]at(11.00,0.5){\parbox{0.9cm}{\centering \S2}};
\fill[AccN!25](14.50,0)rectangle(19.50,1);\draw[white,thick](14.50,0)--(14.50,1);
\node[font=\tiny,rotate=90,anchor=center]at(17.00,0.5){\parbox{0.9cm}{\centering \S3}};
\fill[NavyN!18](19.50,0)rectangle(23.00,1);\draw[white,thick](19.50,0)--(19.50,1);
\node[font=\tiny,rotate=90,anchor=center]at(21.25,0.5){\parbox{0.9cm}{\centering \S4}};
\fill[AccN!25](23.00,0)rectangle(25.50,1);\draw[white,thick](23.00,0)--(23.00,1);
\node[font=\tiny,rotate=90,anchor=center]at(24.25,0.5){\parbox{0.9cm}{\centering \S5}};
\fill[NavyN!18](25.50,0)rectangle(28.00,1);\draw[white,thick](25.50,0)--(25.50,1);
\node[font=\tiny,rotate=90,anchor=center]at(26.75,0.5){\parbox{0.9cm}{\centering \S6}};
\fill[AccN!25](28.00,0)rectangle(29.00,1);\draw[white,thick](28.00,0)--(28.00,1);
\fill[NavyN!18](29.00,0)rectangle(31.50,1);\draw[white,thick](29.00,0)--(29.00,1);
\node[font=\tiny,rotate=90,anchor=center]at(30.25,0.5){\parbox{0.9cm}{\centering \S8}};
\fill[AccN!25](31.50,0)rectangle(32.50,1);\draw[white,thick](31.50,0)--(31.50,1);
\fill[NavyN!18](32.50,0)rectangle(35.00,1);\draw[white,thick](32.50,0)--(32.50,1);
\node[font=\tiny,rotate=90,anchor=center]at(33.75,0.5){\parbox{0.9cm}{\centering \S10}};
\fill[RedN!30](35,0)rectangle(45,1);\node[font=\scriptsize\bfseries,align=center]at(40,0.5){Auto-quiz\\15 Q $\times$ 40 s};
\draw[NavyN](0,0)rectangle(45,1);
\foreach \tt in {0,5,10,15,20,25,30,35,40,45}{\draw[NavyN](\tt,0)--(\tt,-0.12);\node[font=\tiny,below]at(\tt,-0.12){\tt};}
\node[font=\scriptsize,below]at(22.5,-0.45){minutes};
\end{tikzpicture}\end{center}

\begin{center}\small\renewcommand{\arraystretch}{1.2}\begin{tabularx}{\textwidth}{@{}l l X r@{}}\toprule\textbf{Clock} & \textbf{Slides} & \textbf{Segment} & \textbf{Min}\\\midrule

0:00--5:00 & --- & \textbf{Opening Q\&A and recap} (see prompts below) & 5\\

5:00--7:30 & 2--3 & 1. Motivation (2 slides) & 2.5\\

7:30--14:30 & 4--9 & 2. k-Means (6 slides) & 7\\

14:30--19:30 & 10--13 & 3. DBSCAN (4 slides) & 5\\

19:30--23:00 & 14--16 & 4. Gaussian Mixtures (3 slides) & 3.5\\

23:00--25:30 & 17--18 & 5. Hierarchical (2 slides) & 2.5\\

25:30--28:00 & 19--20 & 6. PCA (2 slides) & 2.5\\

28:00--29:00 & 21 & 7. t-SNE \& UMAP (1 slide) & 1\\

29:00--31:30 & 22--23 & 8. Validation (2 slides) & 2.5\\

31:30--32:30 & 24 & 9. Across Disciplines (1 slide) & 1\\

32:30--35:00 & 25--26 & 10. Summary (2 slides) & 2.5\\

35:00--45:00 & --- & \textbf{Auto-quiz} (instructions 1 min, 15 questions $\times$ 40 s = 10 min of answering; collect scores) & 10\\\bottomrule\end{tabularx}\end{center}

Slide 1 is the title slide and is shown while students settle. The teaching segment covers 25 content slides in 30 minutes, an average pace of 0.8 slides per minute. 

Minute allocations are proportional to slide counts; shift time towards the sections students find hardest.

\begin{instr}{Opening Q\&A prompts (first 5 minutes)}\begin{enumerate}\item From Lecture~3: what does the Kalman gain do when the sensor is very precise?\item Why can $R^2$ be negative?\item How does lasso select features?\end{enumerate}Take two or three answers per prompt; correct misconceptions briefly and note any topic that needs revisiting.\end{instr}

\needspace{12\baselineskip}\section{Slide-by-Slide Teaching Guide}

Each subsection gives the slide range, the suggested time, what to emphasise, and the slide titles in order. Numbers refer to pages of the Beamer PDF.

\needspace{8\baselineskip}\subsection{1. Motivation \normalfont\small\color{GreyN}(slides 2--3, about 2.5 min)}

\textbf{Teaching note.} Why discovery matters: you cannot label what you do not know exists.

\begin{center}\small\renewcommand{\arraystretch}{1.1}\begin{tabularx}{\textwidth}{@{}r X@{}}2 & Lecture 4 --- The Core Challenge: Discover What You Don't Know to Look For\\
3 & Case Study Dataset --- OPFOR Archetype Discovery (JANUS)\\\end{tabularx}\end{center}

\needspace{8\baselineskip}\subsection{2. k-Means \normalfont\small\color{GreyN}(slides 4--9, about 7 min)}

\textbf{Teaching note.} Lloyd's algorithm, convergence, assumptions (spherical, similar-sized clusters), elbow at $k=5$, silhouette 0.712, the five archetypes, per-cluster validation and expert agreement.

\begin{center}\small\renewcommand{\arraystretch}{1.1}\begin{tabularx}{\textwidth}{@{}r X@{}}4 & $k$-Means Clustering --- Minimise Within-Cluster Variance (Lloyd's Algorithm)\\
5 & $k$-Means --- Key Properties, Assumptions and Cost\\
6 & Choosing $k$ --- the Elbow Method (WCSS vs $k$)\\
7 & Choosing $k$ --- Silhouette Analysis Confirms $k=5$\\
8 & $k=5$ Solution --- Five OPFOR Manoeuvre Archetypes Discovered\\
9 & $k=5$ Validation --- Per-Cluster Silhouette and Expert Agreement\\\end{tabularx}\end{center}

\needspace{8\baselineskip}\subsection{3. DBSCAN \normalfont\small\color{GreyN}(slides 10--13, about 5 min)}

\textbf{Teaching note.} Density reachability, the 142 noise trajectories interpreted as deception manoeuvres, the $k$-distance plot, MinPts 8 versus 15 (142 versus 318 noise points), and the new doctrine category.

\begin{center}\small\renewcommand{\arraystretch}{1.1}\begin{tabularx}{\textwidth}{@{}r X@{}}10 & DBSCAN --- Density-Based Clustering: Any Shape, and Outliers Flagged as Noise\\
11 & DBSCAN vs $k$-Means --- and the 142 Deception Manoeuvres\\
12 & DBSCAN --- Tuning $\varepsilon$ and MinPts with the $k$-Distance Plot\\
13 & DBSCAN in Operations --- From Noise Points to a New OPFOR Doctrine Category\\\end{tabularx}\end{center}

\needspace{8\baselineskip}\subsection{4. Gaussian Mixtures \normalfont\small\color{GreyN}(slides 14--16, about 3.5 min)}

\textbf{Teaching note.} Soft assignment, AIC/BIC selection, and the E- and M-steps.

\begin{center}\small\renewcommand{\arraystretch}{1.1}\begin{tabularx}{\textwidth}{@{}r X@{}}14 & Gaussian Mixture Models --- Each Cluster Is a Gaussian: Soft Assignment by Probability\\
15 & GMM Case Study --- Model Selection with AIC and BIC\\
16 & The EM Algorithm --- Maximum Likelihood for GMM by Alternating E- and M-Steps\\\end{tabularx}\end{center}

\needspace{8\baselineskip}\subsection{5. Hierarchical \normalfont\small\color{GreyN}(slides 17--18, about 2.5 min)}

\textbf{Teaching note.} Agglomerative clustering, Ward linkage and the dendrogram; ARI 0.847 with $k$-Means.

\begin{center}\small\renewcommand{\arraystretch}{1.1}\begin{tabularx}{\textwidth}{@{}r X@{}}17 & Hierarchical Clustering --- Dendrograms for Analyst Visualisation\\
18 & Ward Dendrogram Case Study --- Two Different Methods, One Answer\\\end{tabularx}\end{center}

\needspace{8\baselineskip}\subsection{6. PCA \normalfont\small\color{GreyN}(slides 19--20, about 2.5 min)}

\textbf{Teaching note.} Eigen-decomposition of the covariance; four components (4.21, 2.84, 1.62, 1.45 of 12) explain 84.4\%; loadings interpretation.

\begin{center}\small\renewcommand{\arraystretch}{1.1}\begin{tabularx}{\textwidth}{@{}r X@{}}19 & Principal Component Analysis --- Directions of Maximum Variance\\
20 & PCA on the 12-Feature Case Study --- Four Components Explain 84.4\%\\\end{tabularx}\end{center}

\needspace{8\baselineskip}\subsection{7. t-SNE \& UMAP \normalfont\small\color{GreyN}(slide 21, about 1 min)}

\textbf{Teaching note.} t-SNE and UMAP for visualisation; only local structure is reliable; never cluster on t-SNE distances.

\begin{center}\small\renewcommand{\arraystretch}{1.1}\begin{tabularx}{\textwidth}{@{}r X@{}}21 & t-SNE \& UMAP --- Non-linear Dimensionality Reduction for Visualisation\\\end{tabularx}\end{center}

\needspace{8\baselineskip}\subsection{8. Validation \normalfont\small\color{GreyN}(slides 22--23, about 2.5 min)}

\textbf{Teaching note.} Internal, external and stability validation; the four-method comparison table.

\begin{center}\small\renewcommand{\arraystretch}{1.1}\begin{tabularx}{\textwidth}{@{}r X@{}}22 & Cluster Validation --- How Do We KNOW the Clusters Are Real?\\
23 & Algorithm Comparison --- Four Clustering Methods on 8,000 OPFOR Trajectories\\\end{tabularx}\end{center}

\needspace{8\baselineskip}\subsection{9. Across Disciplines \normalfont\small\color{GreyN}(slide 24, about 1 min)}

\textbf{Teaching note.} Clustering in biology (cell types), psychology (personality factors), engineering (fault modes) and aeronautics (flight regimes).

\begin{center}\small\renewcommand{\arraystretch}{1.1}\begin{tabularx}{\textwidth}{@{}r X@{}}24 & Clustering \& Dimensionality Reduction Across the Disciplines\\\end{tabularx}\end{center}

\needspace{8\baselineskip}\subsection{10. Summary \normalfont\small\color{GreyN}(slides 25--26, about 2.5 min)}

\textbf{Teaching note.} Summary and look-ahead to anomaly detection.

\begin{center}\small\renewcommand{\arraystretch}{1.1}\begin{tabularx}{\textwidth}{@{}r X@{}}25 & Lecture 4 Summary --- Four Clustering Methods and Three Dimensionality-Reduction Techniques\\
26 & Looking Ahead --- Lecture 5: Unsupervised Learning Part 2 --- Anomaly Detection\\\end{tabularx}\end{center}

\needspace{12\baselineskip}\section{Key Messages}Students should leave the room able to state these in their own words:\begin{keypoint}\begin{enumerate}\item Unsupervised learning finds candidate structure; humans decide whether it is real and meaningful.\item Different algorithms encode different notions of a cluster (centroid, density, probability, hierarchy).\item Agreement between independent methods (ARI) plus expert review is the best evidence a cluster is real.\item Noise points can be the most valuable output: deception and novelty live there.\item Dimensionality reduction is for understanding and pre-processing, not proof.\end{enumerate}\end{keypoint}

\needspace{12\baselineskip}\section{Facilitation Tips}\begin{itemize}\item Run two iterations of $k$-Means on the board with six 1-D points; students then trust the algorithm.\item Show the DBSCAN noise slide as a ``plot twist'': the points $k$-Means forced into clusters were the interesting ones.\item Warn against reading t-SNE cluster sizes or distances: show the same data with two perplexities if you can.\end{itemize}

\needspace{12\baselineskip}\section{Anticipated Questions and Model Answers}

\begin{question}{anticipated question}\textbf{Q:} How do I know $k=5$ is right and not 4 or 6?\par\smallskip\textbf{A:} Elbow, silhouette, BIC and stability across resamples should agree; then experts must be able to name and use the clusters.\end{question}

\begin{question}{anticipated question}\textbf{Q:} Why not always use DBSCAN?\par\smallskip\textbf{A:} It struggles when densities vary and in high dimensions, and $\varepsilon$ is hard to set; $k$-Means is fast and interpretable.\end{question}

\begin{question}{anticipated question}\textbf{Q:} Is PCA required before clustering?\par\smallskip\textbf{A:} Not required, but it reduces noise and computation; keep enough components (here 84\%) and check that clusters survive.\end{question}

\needspace{12\baselineskip}\section{Common Misconceptions}\begin{center}\small\renewcommand{\arraystretch}{1.35}\begin{tabularx}{\textwidth}{@{}>{\raggedright\arraybackslash}X >{\raggedright\arraybackslash}X@{}}\toprule\textbf{Misconception} & \textbf{Correction}\\\midrule``Lowest WCSS gives the best $k$.'' & WCSS always falls as $k$ rises; look for the elbow and confirm.\\
``Clusters found are natural kinds.'' & Algorithms always return clusters, even on noise; validate.\\
``t-SNE shows true distances.'' & Only neighbourhoods are preserved.\\\bottomrule\end{tabularx}\end{center}

\needspace{12\baselineskip}\section{In-Class Activities and Discussion Prompts}\label{sec:act}\begin{itemize}\item \textbf{Name the archetype:} show cluster centroids without labels; teams propose doctrinal names and the evidence they would need to confirm them.\item \textbf{Silhouette by hand} for one point with $a=2$, $b=5$: $s=0.6$.\end{itemize}

\needspace{12\baselineskip}\section{Run It Live: Studio Demonstrations}\label{sec:live}

Project \texttt{Lecture4\_Clustering\_Lab.html} (it is a single HTML file and runs offline in any browser; every result updates as an input changes). Each demonstration below follows the same rhythm: \textbf{set} the inputs, \textbf{ask} the class to commit to a prediction (a show of hands or a one-word answer), then \textbf{reveal} the result and \textbf{land} the point. The numbers are exactly what the page shows for the settings given; simulations use fixed random seeds, so they repeat every time. The 45-minute plan has little slack: choose the one demonstration marked \emph{first choice} for class, and use the others in a tutorial, a second session, or as homework with the drill cards at the bottom of the page.

\begin{live}{Demo 1: When k-means is the wrong tool (about 4 min, with \S3 DBSCAN) --- \emph{first choice}}\textbf{Set.} Demo~1: dataset \textsf{Two curved groups}, $k=2$, press \textsf{Run to convergence}.\par\smallskip\textbf{Ask.} ``Two groups, two clusters. Will k-means find them?''\par\smallskip\textbf{Reveal.} WCSS 303.6, silhouette 0.507, but agreement with the true groups (ARI) only \textbf{0.42}: the boundary cuts across both curves. Demo~2 even says silhouette peaks at $k=8$. Demo~3 (DBSCAN, $\varepsilon=0.9$, minPts 5) finds 2 clusters with \textbf{ARI 1.000}, against 0.375 for the best k-means.\par\smallskip\textbf{Land the point.} k-means assumes compact, convex groups, and so do the elbow and silhouette. Look at the clusters on the map before trusting any index.\end{live}

\begin{live}{Demo 2: One $\varepsilon$ for every density (about 2 min, with \S3 DBSCAN)}\textbf{Set.} Demo~3 on the curved groups.\par\smallskip\textbf{Ask.} ``What happens if $\varepsilon$ is halved? Doubled?''\par\smallskip\textbf{Reveal.} $\varepsilon=0.5$: 8 clusters, 9 noise points, ARI 0.281. $\varepsilon=1.3$: everything merges into 1 cluster, ARI 0.000.\par\smallskip\textbf{Land the point.} DBSCAN's results hinge on $\varepsilon$. Use a $k$-distance plot to choose it, and HDBSCAN when densities differ.\end{live}

\begin{live}{Demo 3: Units hijack PCA (about 2 min, with \S6 PCA)}\textbf{Set.} Demo~4: correlation 0.3, feature~2 recorded in units 4$\times$ larger, preprocessing \textsf{Raw}.\par\smallskip\textbf{Ask.} ``PC1 explains 95\% of the variance. Have we found strong structure?''\par\smallskip\textbf{Reveal.} Raw: PC1 explains \textbf{95.4\%} at 85.1$^\circ$, almost exactly along feature~2. Switch to \textsf{Standardised}: \textbf{67.3\%} at 45.0$^\circ$.\par\smallskip\textbf{Land the point.} Before standardising, PCA mostly reflects measurement units. Standardise unless the units are genuinely comparable.\end{live}

\needspace{12\baselineskip}\section{Consistency Notes on the Source Material}These notes record where the original PowerPoint deck contained an inconsistency and how the current slides, reference text, quiz and assignment handle it. Students who have seen the old deck may raise them.\begin{itemize}\item In the source deck the listed cluster sizes sum to 6,000 rather than 8,000, and the GMM mixture weights sum to 0.75 rather than 1. They appear to describe a 6,000-trajectory subset; the reference text says so. Use it as a ``check that partitions account for the whole dataset'' moment.\end{itemize}

\needspace{12\baselineskip}\section{Quiz Administration and Answer Key}

\begin{sayit}``You will now take a 15-question quiz. Each question has four options and exactly one correct answer. You have 40 seconds per question; the quiz advances automatically when time runs out. A correct answer scores $+1$, a wrong answer $-0.25$, and an unanswered question 0. The maximum is 15 and the minimum is $-3.75$. Because of the negative marking, guessing blindly among four options gains nothing on average --- answer when you can eliminate options. Your score and a full review appear at the end.''\end{sayit}

\textbf{Delivery.} Open \texttt{Lecture4\_Quiz.html} on each student device (or project it and collect answers on paper). The page enforces the timer and computes the score with negative marking. Release \texttt{Lecture4\_Answers\_Assignment.html} only after all students have finished.

\subsection*{Answer key}\begin{center}\small\renewcommand{\arraystretch}{1.3}\begin{tabular}{|c|c|c|c|c|c|c|c|c|c|c|c|c|c|c|}\hline \textbf{1} & \textbf{2} & \textbf{3} & \textbf{4} & \textbf{5} & \textbf{6} & \textbf{7} & \textbf{8} & \textbf{9} & \textbf{10} & \textbf{11} & \textbf{12} & \textbf{13} & \textbf{14} & \textbf{15}\\\hline A & D & C & C & A & B & C & A & D & D & C & A & B & C & D\\\hline\end{tabular}\end{center}

\textbf{Interpreting scores.} 12 or more: secure command of the lecture. 9--11.75: good, review missed items. Below 9: revisit the reference text before the assignment. If more than half the class misses a question, revisit that point in the next lecture's opening Q\&A.

\subsection*{Questions, options and rationale}

{\small

\noindent\textbf{Q1.} k-Means minimises:\begin{itemize}[leftmargin=2.2em,itemsep=0pt,topsep=1pt]\item[\textbf{A.}] \textbf{The within-cluster sum of squared distances to centroids (WCSS)}~{\color{GrnN}$\checkmark$}\item[B.] The number of clusters\item[C.] The silhouette\item[D.] The likelihood of a Gaussian mixture\end{itemize}\noindent{\color{GreyN}\textit{Rationale:}} $J=\sum_j\sum_{x\in C_j}\|x-\mu_j\|^2$.\par\medskip

\noindent\textbf{Q2.} Why does Lloyd's algorithm always converge?\begin{itemize}[leftmargin=2.2em,itemsep=0pt,topsep=1pt]\item[A.] It finds the global optimum\item[B.] k is fixed\item[C.] It uses gradient descent with momentum\item[\textbf{D.}] \textbf{Neither the assignment nor the update step can increase J, which is bounded below}~{\color{GrnN}$\checkmark$}\end{itemize}\noindent{\color{GreyN}\textit{Rationale:}} Monotone decrease to a local minimum; restarts and k-means++ reduce poor minima.\par\medskip

\noindent\textbf{Q3.} WCSS keeps decreasing as k increases. Therefore:\begin{itemize}[leftmargin=2.2em,itemsep=0pt,topsep=1pt]\item[A.] Choose the largest k\item[B.] WCSS is useless\item[\textbf{C.}] \textbf{Choose k at the elbow and confirm with silhouette/domain knowledge}~{\color{GrnN}$\checkmark$}\item[D.] Always choose k = 2\end{itemize}\noindent{\color{GreyN}\textit{Rationale:}} WCSS is zero when k = N; look for diminishing returns (k = 5 here).\par\medskip

\noindent\textbf{Q4.} A point has silhouette $s(i)=-0.2$. This means:\begin{itemize}[leftmargin=2.2em,itemsep=0pt,topsep=1pt]\item[A.] It is well clustered\item[B.] It is noise\item[\textbf{C.}] \textbf{It is closer on average to another cluster than to its own: probably misassigned}~{\color{GrnN}$\checkmark$}\item[D.] The clustering is perfect\end{itemize}\noindent{\color{GreyN}\textit{Rationale:}} $s<0\iff b(i)<a(i)$.\par\medskip

\noindent\textbf{Q5.} The INFILTRATION archetype was significant because:\begin{itemize}[leftmargin=2.2em,itemsep=0pt,topsep=1pt]\item[\textbf{A.}] \textbf{It was not in the analysts' label scheme: a tactic discovered without labels}~{\color{GrnN}$\checkmark$}\item[B.] It was the largest cluster\item[C.] It had the highest speed\item[D.] DBSCAN labelled it noise\end{itemize}\noindent{\color{GreyN}\textit{Rationale:}} Unsupervised learning surfaces categories nobody had defined; experts then validated it.\par\medskip

\noindent\textbf{Q6.} In DBSCAN, a border point is:\begin{itemize}[leftmargin=2.2em,itemsep=0pt,topsep=1pt]\item[A.] A point with at least MinPts neighbours\item[\textbf{B.}] \textbf{A non-core point within epsilon of a core point}~{\color{GrnN}$\checkmark$}\item[C.] An outlier\item[D.] A centroid\end{itemize}\noindent{\color{GreyN}\textit{Rationale:}} Core: at least MinPts within epsilon; border: reachable from a core point; noise: neither.\par\medskip

\noindent\textbf{Q7.} DBSCAN's 142 noise trajectories were interpreted as:\begin{itemize}[leftmargin=2.2em,itemsep=0pt,topsep=1pt]\item[A.] Sensor errors to discard\item[B.] Patrols\item[\textbf{C.}] \textbf{Deliberate deception manoeuvres}~{\color{GrnN}$\checkmark$}\item[D.] Duplicate records\end{itemize}\noindent{\color{GreyN}\textit{Rationale:}} In defence, 'noise' is often the most interesting signal; k-Means would have hidden them.\par\medskip

\noindent\textbf{Q8.} Increasing MinPts from 8 to 15 (same epsilon) will:\begin{itemize}[leftmargin=2.2em,itemsep=0pt,topsep=1pt]\item[\textbf{A.}] \textbf{Increase the number of noise points (318 vs 142)}~{\color{GrnN}$\checkmark$}\item[B.] Reduce the number of noise points\item[C.] Have no effect\item[D.] Merge all clusters\end{itemize}\noindent{\color{GreyN}\textit{Rationale:}} A stricter density requirement makes more points non-core.\par\medskip

\noindent\textbf{Q9.} GMM responsibilities $r_{ik}$ are:\begin{itemize}[leftmargin=2.2em,itemsep=0pt,topsep=1pt]\item[A.] Hard 0/1 assignments\item[B.] Distances\item[C.] Cluster sizes\item[\textbf{D.}] \textbf{Posterior probabilities that point i came from component k}~{\color{GrnN}$\checkmark$}\end{itemize}\noindent{\color{GreyN}\textit{Rationale:}} Soft membership: e.g. T481 had 0.62 assault / 0.38 infiltration.\par\medskip

\noindent\textbf{Q10.} For selecting the number of GMM components you should use:\begin{itemize}[leftmargin=2.2em,itemsep=0pt,topsep=1pt]\item[A.] The raw log-likelihood\item[B.] Training time\item[C.] The number of features\item[\textbf{D.}] \textbf{AIC or BIC, which penalise extra parameters}~{\color{GrnN}$\checkmark$}\end{itemize}\noindent{\color{GreyN}\textit{Rationale:}} Likelihood always rises with K; AIC/BIC both chose K = 5.\par\medskip

\noindent\textbf{Q11.} Ward linkage merges the pair of clusters that:\begin{itemize}[leftmargin=2.2em,itemsep=0pt,topsep=1pt]\item[A.] Are farthest apart\item[B.] Have the most points\item[\textbf{C.}] \textbf{Least increases the total within-cluster variance}~{\color{GrnN}$\checkmark$}\item[D.] Are nearest by single linkage\end{itemize}\noindent{\color{GreyN}\textit{Rationale:}} It produces compact clusters and suits defence analysis.\par\medskip

\noindent\textbf{Q12.} ARI = 0.847 between Ward and k-Means means:\begin{itemize}[leftmargin=2.2em,itemsep=0pt,topsep=1pt]\item[\textbf{A.}] \textbf{Two different methods produced very similar partitions: evidence the clusters are real}~{\color{GrnN}$\checkmark$}\item[B.] The methods disagree\item[C.] The data are noise\item[D.] k is wrong\end{itemize}\noindent{\color{GreyN}\textit{Rationale:}} ARI is 1 for identical partitions and about 0 for chance agreement.\par\medskip

\noindent\textbf{Q13.} PCA's first principal component is:\begin{itemize}[leftmargin=2.2em,itemsep=0pt,topsep=1pt]\item[A.] The feature with the largest mean\item[\textbf{B.}] \textbf{The eigenvector of the covariance matrix with the largest eigenvalue}~{\color{GrnN}$\checkmark$}\item[C.] A random direction\item[D.] The cluster centroid\end{itemize}\noindent{\color{GreyN}\textit{Rationale:}} It maximises projected variance; its eigenvalue is that variance.\par\medskip

\noindent\textbf{Q14.} Four PCs with eigenvalues 4.21, 2.84, 1.62, 1.45 out of a total of 12 explain about:\begin{itemize}[leftmargin=2.2em,itemsep=0pt,topsep=1pt]\item[A.] 35\%\item[B.] 59\%\item[\textbf{C.}] \textbf{84\%}~{\color{GrnN}$\checkmark$}\item[D.] 100\%\end{itemize}\noindent{\color{GreyN}\textit{Rationale:}} $(4.21+2.84+1.62+1.45)/12=10.12/12\approx0.84$.\par\medskip

\noindent\textbf{Q15.} In a t-SNE plot, which statement is correct?\begin{itemize}[leftmargin=2.2em,itemsep=0pt,topsep=1pt]\item[A.] Distances between clusters are meaningful\item[B.] t-SNE is linear\item[C.] Cluster sizes show true densities\item[\textbf{D.}] \textbf{Only local neighbourhood structure is reliably preserved}~{\color{GrnN}$\checkmark$}\end{itemize}\noindent{\color{GreyN}\textit{Rationale:}} Never interpret between-cluster distances or sizes in t-SNE; UMAP preserves somewhat more global structure.\par\medskip

}

\needspace{12\baselineskip}\section{Assignment: Marking Scheme and Model Solutions}

\textbf{Instructions.} Five problems, 100 marks, plus a 20-mark bonus problem (Problem 6). Show all working. Reference: Lecture 4 reference text.

\begin{center}\small\renewcommand{\arraystretch}{1.25}\begin{tabularx}{0.85\textwidth}{@{}l X r@{}}\toprule\textbf{\#} & \textbf{Problem} & \textbf{Marks}\\\midrule1 & k-Means by hand & 20\\
2 & Silhouette values & 15\\
3 & One EM iteration & 25\\
4 & PCA by hand & 20\\
5 & DBSCAN by hand & 20\\
6 & Problem 6 (bonus): Discovering emitter families & 20\\\midrule & \textbf{Total} & \textbf{100} + 20 bonus\\\bottomrule\end{tabularx}\end{center}

\begin{instr}{General marking guidance}\begin{itemize}
\item Marks for each part are shown in square brackets in the problem statement.
\item Award method marks for correct working even when an arithmetic slip carries through; deduct once for the slip, not again for its consequences.
\item For interpretive parts, full marks need (i) the correct conclusion, (ii) the reason, and (iii) an operational implication or caveat.
\item Answers that state a number without uncertainty, baseline or caveat where one is required lose a mark per part.
\item Suggested deadline: one week after the lecture; return marked work before the next quiz on the same topic.
\end{itemize}\end{instr}

\needspace{8\baselineskip}\subsection{Problem 1: k-Means by hand \normalfont\small\color{GreyN}(20 marks)}

\begin{prob}One-dimensional speeds: 1, 2, 3, 10, 11, 12, 20. Run $k$-Means with $k=2$ and initial centroids 1 and 20.

(a) Perform iterations until convergence, showing assignments and centroids. [10]\newline (b) Compute the final WCSS. [5]\newline (c) What would $k=3$ give, and how would you decide between 2 and 3? [5]\end{prob}

\begin{soln}\textbf{(a)} Iteration 1: $\{1,2,3,10\}$ (10 is 9 from 1 and 10 from 20) and $\{11,12,20\}$; centroids 4 and 14.33. Iteration 2: 10 is now 6 from 4 and 4.33 from 14.33, so it moves: $\{1,2,3\}$, $\{10,11,12,20\}$; centroids 2 and 13.25. Iteration 3: no change $\rightarrow$ converged.

\textbf{(b)} Cluster 1: $1+0+1=2$. Cluster 2: $3.25^2+2.25^2+1.25^2+6.75^2=10.56+5.06+1.56+45.56=62.75$. $J=\mathbf{64.75}$.

\textbf{(c)} $k=3$ gives $\{1,2,3\}$, $\{10,11,12\}$, $\{20\}$ with $J=2+2+0=4$, a huge drop. Since WCSS always falls with $k$, judge by the size of the drop, silhouette, and meaning: a singleton cluster may be an outlier (DBSCAN would call it noise) rather than a real group.\end{soln}

\needspace{8\baselineskip}\subsection{Problem 2: Silhouette values \normalfont\small\color{GreyN}(15 marks)}

\begin{prob}Point P is in cluster A; its distances to the other members of A are 1 and 2; to the members of the nearest other cluster B they are 4, 5 and 6. Point Q in cluster A has mean distance 3 to A and 2.5 to B.

(a) Compute $s(P)$ and $s(Q)$. [8]\newline (b) Interpret both and state what you would do about Q. [7]\end{prob}

\begin{soln}\textbf{(a)} $a(P)=1.5$, $b(P)=5$, $s(P)=(5-1.5)/5=\mathbf{0.70}$. $s(Q)=(2.5-3)/3=\mathbf{-0.17}$.

\textbf{(b)} P is well clustered. Q is closer on average to B than to its own cluster: probably misassigned, or a boundary/ambiguous case. Reassign it, examine it as a possible anomaly, or use a soft method (GMM) to report its membership probabilities.\end{soln}

\needspace{8\baselineskip}\subsection{Problem 3: One EM iteration \normalfont\small\color{GreyN}(25 marks)}

\begin{prob}A two-component 1-D GMM has $\pi_1=\pi_2=0.5$, $\mu_1=0$, $\mu_2=4$, $\sigma_1=\sigma_2=1$. Data: $x=0, 1.5, 4$.

(a) E-step: compute the responsibilities for all three points. [12]\newline (b) M-step: update $\pi_1$ and $\mu_1$. [8]\newline (c) Explain why the likelihood cannot decrease, and name two practical pitfalls of EM. [5]\end{prob}

\begin{soln}\textbf{(a)} $\mathcal N(x;\mu,1)=0.3989e^{-(x-\mu)^2/2}$. $x=0$: $0.3989$ vs $0.3989e^{-8}=1.34\times10^{-4}$; $r_1=\mathbf{0.9997}$. $x=1.5$: $0.3989e^{-1.125}=0.1295$ vs $0.3989e^{-3.125}=0.0175$; $r_1=0.1295/0.1470=\mathbf{0.881}$, $r_2=0.119$. $x=4$: symmetric to $x=0$: $r_1=\mathbf{0.0003}$.

\textbf{(b)} $N_1=0.9997+0.881+0.0003=1.881$; $\pi_1=1.881/3=\mathbf{0.627}$; $\mu_1=(0.9997\cdot0+0.881\cdot1.5+0.0003\cdot4)/1.881=1.323/1.881=\mathbf{0.703}$.

\textbf{(c)} EM maximises a lower bound on the log-likelihood that touches it at the current parameters (E-step) and then raises the bound (M-step), so the likelihood never falls. Pitfalls: local maxima (use multiple starts), collapsing components with $\sigma\to0$ (regularise), label switching, and sensitivity to initialisation.\end{soln}

\needspace{8\baselineskip}\subsection{Problem 4: PCA by hand \normalfont\small\color{GreyN}(20 marks)}

\begin{prob}Two standardised features (speed, direction change) have covariance matrix $C=\begin{pmatrix}3&1\\1&3\end{pmatrix}$.

(a) Find the eigenvalues and unit eigenvectors. [10]\newline (b) Compute each component's explained-variance ratio. [4]\newline (c) Project the point $(2,0)$ onto PC1 and interpret PC1. [6]\end{prob}

\begin{soln}\textbf{(a)} $\det(C-\lambda I)=(3-\lambda)^2-1=0\Rightarrow\lambda=\mathbf{4},\mathbf{2}$. For $\lambda=4$: $(1,1)/\sqrt2$; for $\lambda=2$: $(1,-1)/\sqrt2$.

\textbf{(b)} EVR $=4/6=\mathbf{66.7\%}$ and $2/6=\mathbf{33.3\%}$.

\textbf{(c)} Score $=(2,0)\cdot(1,1)/\sqrt2=\mathbf{1.41}$. PC1 weights speed and direction change equally: an overall 'activity' axis on which fast, manoeuvring trajectories score high; PC2 contrasts the two (fast but straight vs slow but erratic).\end{soln}

\needspace{8\baselineskip}\subsection{Problem 5: DBSCAN by hand \normalfont\small\color{GreyN}(20 marks)}

\begin{prob}Points: 1, 2, 3, 4, 10, 11, 12, 20 with $\varepsilon=1.5$ and MinPts = 3 (a point counts itself).

(a) Label each point core, border or noise and give the clusters. [10]\newline (b) Repeat with MinPts = 4. [5]\newline (c) Explain why the noise point might matter more than either cluster in an intelligence context. [5]\end{prob}

\begin{soln}\textbf{(a)} Neighbourhoods: 1: $\{1,2\}$; 2: $\{1,2,3\}$; 3: $\{2,3,4\}$; 4: $\{3,4\}$; 10: $\{10,11\}$; 11: $\{10,11,12\}$; 12: $\{11,12\}$; 20: $\{20\}$. Core: 2, 3, 11. Border: 1, 4 (near 2 and 3), 10, 12 (near 11). Noise: 20. Clusters $\{1,2,3,4\}$ and $\{10,11,12\}$.

\textbf{(b)} No point has 4 neighbours within 1.5, so there are no core points and everything is noise: MinPts too large for this density.

\textbf{(c)} An isolated behaviour unlike any routine pattern can be a new tactic, deception or a sensor fault. DBSCAN surfaces it explicitly for analysts, whereas $k$-Means would absorb it into the nearest group (as with the 142 deception manoeuvres).\end{soln}

\needspace{8\baselineskip}\subsection{Problem 6 (bonus): Discovering emitter families \normalfont\small\color{GreyN}(20 marks)}

\begin{prob}Nine radar intercepts are described by pulse repetition frequency (PRF, 800--3,000 Hz) and pulse width (0.5--5 $\mu$s). After standardisation their coordinates are: P1 $(0,0)$, P2 $(0.4,0.2)$, P3 $(0.2,0.5)$, P4 $(0.5,0.6)$, P5 $(3,3)$, P6 $(3.3,2.8)$, P7 $(3.1,3.4)$, P8 $(1.6,1.7)$, P9 $(5,0)$.

(a) Why must the raw features be standardised before clustering? [3]\newline (b) Run DBSCAN with $\varepsilon=0.6$ and MinPts $=3$ (counting the point itself). Label every point core, border or noise and give the clusters. [7]\newline (c) Compute the silhouette of P4, taking clusters A $=\{$P1--P4$\}$ and B $=\{$P5--P7$\}$. [4]\newline (d) Along the PRF axis alone, a two-component mixture has family A $\sim\mathcal N(1.2,0.3^2)$ with weight 0.6 and family B $\sim\mathcal N(2.4,0.4^2)$ with weight 0.4 (standardised units). Compute the responsibility of family A for an intercept at 1.8, and say what an analyst should do with it. [6]\end{prob}

\begin{soln}\textbf{(a)} Raw PRF differences run to thousands of hertz while pulse-width differences are a few microseconds, so Euclidean distance would be decided almost entirely by PRF. Standardising each feature to zero mean and unit variance gives both a fair say. Record the means and standard deviations, because new intercepts must be transformed the same way.

\textbf{(b)} Neighbour counts within 0.6: P1 3, P2 4, P3 4, P4 3, P5 3 $\Rightarrow$ \textbf{core}. P6 2 and P7 2, but each lies within 0.6 of the core point P5 $\Rightarrow$ \textbf{border} (P6--P7 is 0.63, just outside). P8 and P9 have only themselves $\Rightarrow$ \textbf{noise}. Clusters: A $=\{$P1, P2, P3, P4$\}$; B $=\{$P5, P6, P7$\}$. P8 lies between the families (mean distance 1.91 to A, 2.07 to B); $k$-means would force it into A, whereas DBSCAN leaves it for an analyst. P9 is isolated: a candidate new emitter type.

\textbf{(c)} $a$ = mean distance from P4 to P1, P2, P3 $=0.50$; $b$ = mean distance to P5, P6, P7 $=3.62$; $s=(3.62-0.50)/3.62=\mathbf{0.86}$: firmly inside its cluster.

\textbf{(d)} $0.6\,\mathcal N(1.8;1.2,0.3^2)=0.108$ and $0.4\,\mathcal N(1.8;2.4,0.4^2)=0.130$, so $r_A=0.108/0.238=\mathbf{0.45}$. The intercept is genuinely ambiguous (an intercept at 2.0 would give 0.09). Report it as uncertain and seek a second feature (pulse width, scan pattern) or a further intercept rather than forcing a label; soft assignments are exactly what a GMM is for.\end{soln}

\end{document}